\section{Outstanding analytical tasks}\label{sec:outstanding}

The construction realizes the canonical family within the supplied ideal
conditional model. The collision analysis describes a hypothetical event.
Section~\ref{sec:highcount} now adds two assessed scalar consequences:
complete high-count localization with both signs, and absolute value/first-
derivative errors for finite count reconstruction. These settle neither
the source-specific collision question nor an advantage over direct source
analysis. They define useful constraints on the next conceptual pass.

\subsection{Identify an independent source-specific consequence}
The event conditions are the two vanished balances in
\eqref{eq:momenttriple} and, at a double event, a nonzero signed second
moment. Source-sector covariance and count-sector transfers describe the
same signed information. Equation~\eqref{eq:scorebalance} is also a
consequence of the existing balances, not an additional nonvanishing
condition. Theorem~\ref{thm:highcount} excludes eventual one-sign literal
count sectors for theta itself at fixed nonzero real $k$; renaming that
premise cannot supply a collision mechanism. The absolute bounds in
Theorem~\ref{thm:counttail} pay omitted sectors but do not create a
relative margin when the readout and its first derivative vanish.

The next question is whether the complete source imposes an independently
proved joint relation or quantitative constraint on the signed balances,
or yields another useful analytical consequence. State the actual source
premise, parameter/multiplicity range and control that tests it. A rigidity
or transport relation can be useful before a scalar bound is selected,
but requiring $C$ and $C_k$ never to vanish together only restates the
target. All global proposals retain the endpoint, confinement and
equivalence qualifications of existing heat-flow criteria.

\subsection{Freeze a bounded task only after a mathematical decision}
TND-D-B02's selected diagnostic and the E03/E04 and E06/E07 proof/assessment
lanes are complete within their stated scopes. The next-session brainstorm
is recorded as TND-D-B03, recommended Astra at effort ultra, ready but
unactivated at this manuscript integration. No new analytical execution
contract is selected here.

A subsequent contract must name one theorem, identity or estimate, the
source information, parameter/readout ranges, derivative orders, complete-
tail method and success or failure test. Compare with direct use of
exactly the same source and uncertainty information. The count labels add
no independent arithmetic data. The generic covariance bound
\eqref{eq:covbound} controls speed without controlling the dangerous signed
balance, and \eqref{eq:variance} gives nonzero observation variance at zero
contrast. A conditioning or practical reconstruction claim therefore
needs an explicit information/error model. The present conservative
threshold is finite but uncomputed and unoptimized; it is not a cost claim.

\subsection{Cover multiplicities and the desired global range}
The double-event picture is local and nondegenerate.
Proposition~\ref{prop:multiple} identifies higher-order data but supplies
neither constants uniform in multiplicity nor a mechanism covering all
events. An exclusion must cover the relevant multiplicities or justify
why omitted ones cannot arise. It must also handle temporal endpoints and
large readout parameters. The compact-real-$k$ phase theorem and the
all-real-$k$ absolute tail theorem have different scopes; neither supplies
an all-height relative collision margin. Fixed finite-order domains do not
imply a uniform graph or energy estimate as the count increases. The
closed apparatus slab gives no infinite-time conclusion.

\subsection{Separate analytical reconstruction from physical access}
The theta loader, spatial imprint, complete angular ensemble, absolute
scale and transported no-count readout remain supplied. Deriving them
intrinsically or specifying a microscopic implementation and its resource
or error bounds is a separate physical task. The nonzero-count quantities
are valid analytical traces, and \eqref{eq:tailreconstruction} is an
analytical finite reconstruction. An experiment using those joint
readouts needs its own explicit access specification. A changed controller,
source, clock or measurement law requires a revised construction before
transferring a claim. Changing $\gamma$ only changes the laboratory clock.

\subsection{Preserve assessment boundaries and use computation as support}
The realization retains its internal conditional assessment. The new
concentration, phase, complete-tail and first-derivative reconstruction
results have separate-worker, same-model internal LLM assessments with
shared premises. Only the specifically consumed count identities and
source-score identity share that acceptance. Other fold, higher-
multiplicity, finite-transfer and whole-channel kernel diagnostics retain
proof-author status; manuscript production does not supply specialist
independence or a wider assessment.

No numerical evidence is needed for these analytical results. Supporting
computation may later evaluate a declared constant, compare conditioning
under a fixed error model or test a finite-domain question, with complete-
tail errors and a distinction between floating and certified evidence.
It cannot replace the independent signed source estimate needed for
collision control. No new numerical campaign is selected by this draft.
